\documentclass[10pt,a4paper]{article}
\usepackage[utf8]{inputenc}
\usepackage[T1]{fontenc}
\usepackage[spanish]{babel}
\usepackage{amsmath,amssymb}
\usepackage[left=2cm,right=2cm,top=1.2cm,bottom=1.1cm]{geometry}
\usepackage{xcolor,array,tabularx,enumitem}
\definecolor{azul}{HTML}{1B3A5C}
\definecolor{gris}{HTML}{555555}
\usepackage[colorlinks=true,urlcolor=azul,linkcolor=azul]{hyperref}
\pagestyle{empty}
\setlength{\parindent}{0pt}
\setlength{\parskip}{3pt}
\setlength{\emergencystretch}{2em}
\newcommand{\seccion}{1}
\newcommand{\duracion}{5 minutos y 15 segundos}
\newcommand{\sha}{\href{https://github.com/CataVergara/EID-Calculo-Intermedio-MAT1188-/tree/ac07f4fa2322818e2686b4b5d0b36dbbf1d10562}{\texttt{ac07f4fa2322}}}
\newcommand{\apartado}[1]{\vspace{4pt}\textbf{\color{azul}#1}\par}
\newcolumntype{E}{>{\raggedright\arraybackslash\bfseries}p{3.3cm}}
\newcolumntype{C}{>{\raggedright\arraybackslash}X}
\renewcommand{\arraystretch}{1.1}
\setlength{\tabcolsep}{5pt}
\hypersetup{pdftitle={Ficha de entrega EID MAT1188, Grupo 6},pdfauthor={Grupo 6}}

\begin{document}
\small
{\Large\bfseries\color{azul} Ficha de entrega: EID Cálculo Intermedio}\\[3pt]
{\color{gris} MAT1188. Ingeniería Civil Informática. Grupo 6. Año 2026.}

\begin{tabularx}{\textwidth}{|E|C|}
\hline
Curso y sección & MAT1188, Cálculo Intermedio. Sección: \seccion \\\hline
Integrantes & Catalina Javiera Vergara Alarcón\newline Ignacio Alexander De Celis Valencia\newline Marcelo Fabián Santana Astorga \\\hline
Tema & Integral definida y área geométrica: contribuciones positivas y negativas \\\hline
Pregunta central & ¿Por qué una integral definida puede ser menor que el área geométrica total? \\\hline
Video & \texttt{MAT1188\_EID\_G06\_2026.mp4}\newline Duración real: \duracion \\\hline
Repositorio & \url{https://github.com/CataVergara/EID-Calculo-Intermedio-MAT1188-} \\\hline
Versión del código & Commit del código utilizado: \sha \\\hline
\end{tabularx}

\apartado{Ejemplo y respuesta}
Usamos $f(x)=x^2-4$ en $[-3,3]$. Es continua y sus raíces son $-2$ y $2$.
Es negativa en $(-2,2)$ y vale cero en las raíces.
Las regiones sobre el eje tienen área total $A^+=14/3$.
La región bajo el eje tiene área $A^-=32/3$.
\[
\int_{-3}^{3}f(x)\,dx=A^+-A^-=-6,\qquad
\int_{-3}^{3}|f(x)|\,dx=A^++A^-=\frac{46}{3}\approx15{,}33.
\]
\textbf{Respuesta:} la integral suma los aportes con su signo. Las regiones bajo el eje aportan valores negativos. El área geométrica suma las áreas de todas las regiones. Por eso, si $A^->0$, la integral es menor que el área total. Si $A^-=0$, ambas coinciden.

\apartado{Referencias}
\begin{itemize}[leftmargin=1.2em,itemsep=1pt,topsep=1pt]
\item Apunte base: \emph{Introducción a la integral}. Cálculo Intermedio, Ingeniería Civil Informática, 2026. Material del curso en Blackboard.
\item Stewart, J. (2012). \emph{Cálculo de una variable. Trascendentes tempranas}. 7.\textordfeminine\ edición. Cengage Learning. Páginas 373 y 378. \href{https://eva.interior.udelar.edu.uy/pluginfile.php/96366/mod_resource/content/1/Stewart.\%20C\%C3\%A1lculo\%20de\%20una\%20variable..pdf}{PDF de la Universidad de la República}.
\item Thomas, G. B., Weir, M. D. y Hass, J. (2010). \emph{Cálculo. Una variable}. 12.\textordfeminine\ edición. Pearson. Páginas 280 y 281. \href{https://fcen.uncuyo.edu.ar/upload/calculo-de-una-variable-thomas-12ed.pdf}{PDF de la Universidad Nacional de Cuyo}.
\item Manim Community. Versión 0.22.0. Documentación: \url{https://docs.manim.community}.
\end{itemize}
Las ediciones y páginas se verificaron en los PDF. Referencias completas, FFmpeg y registro de IA en \href{https://github.com/CataVergara/EID-Calculo-Intermedio-MAT1188-/blob/main/docs/REFERENCIAS.md}{el registro del repositorio}.

\apartado{Declaración de uso de inteligencia artificial}
\begin{tabularx}{\textwidth}{|E|C|}
\hline
Herramientas & Claude, de Anthropic, Codex y Whisper, de OpenAI. \\\hline
Apoyo recibido & Claude: guía de trabajo, propuesta del ejemplo, primer borrador del código y correcciones de texto y distribución visual.\newline
Codex: revisión de código, escenas, cálculos, guion, ficha y referencias. Preparación del montaje con FFmpeg.\newline Whisper: transcripción local para elegir tomas y localizar repeticiones. \\\hline
Comprobación & Se compararon los resultados con las pruebas del repositorio y con $F(x)=x^3/3-4x$. Se revisaron las escenas, los tiempos y los recortes. El equipo debe comprobar esta declaración antes de entregar. \\\hline
Voces y revisión final & El montaje usa las grabaciones de Catalina, Ignacio y Marcelo. No se generaron ni clonaron voces. Se recortaron repeticiones sin cambiar la velocidad. El equipo debe escuchar el video completo antes de entregarlo. \\\hline
\end{tabularx}

\end{document}
